% AI-written, AI cross-reviewed; no human sign-off; not a claim that RH is proved
% RH Notes blueprint source excerpts for Z1.1.04
% Verbatim excerpts in the notation of the parent blueprint. The labelled editions and source ranges appear below.
% See the unit YAML for the current proof status. Earlier-round passages retain their original wording.
\section*{Z1.1.04. Rank-three seeds at prescribed weights}

% Parent source SHA-256: 1af64922b98ed2310744411b1f18b25cc426c97220646321471836e1a3a409d6
% blueprint/ch_P.tex lines 12441 through 12471; r5:Z1.1.04
\subsection{Z1.1.04. Building the rank-three seed in a specified weight class}\label{r5:Z1.1.04}
Assume $\mathrm{PropEGW}$, $\mathrm{SeedInW}$ and $\mathrm{AdAutW}$. Fix the field pair and characters of $\mathrm{FieldPairCharacterTripleConditions}$ with $\chi'\ne\eta$, the nonzero level $M$ divisible by $2\ell$, its unramified splitting and the stated inertia condition on $\chi'$. For every bound $B$ we must produce $m\ge B$ congruent to $j$ modulo $\ell-1$ and an openly irreducible $r$ over $L$, with weights $(m,0,-m)$, residual characters $(\Psi,1,\Psi^{-1})$ for $\Psi=\epsilon^j\chi'$, away characters $(\chi',1,(\chi')^{-1})$ restricted to $L$, and automorphy data for $r\otimes\epsilon^{-1}$. Here $\epsilon$ is cyclotomic; an open-irreducibility assertion means irreducibility after restriction to every open subgroup.

{\footnotesize \nolinkurl{UpperTheory.V.Z.Sym.UnitaryThreeAutomorphyDense} (\nolinkurl{Development/AutomorphicArithmetic/AutomorphicVarieties/ModularityArgument/ThresholdArguments/DenseThresholdCases/SymUse.lean}\allowbreak:27--34). \nolinkurl{UpperTheory.V.Sd.IsSeedAlt} (\nolinkurl{Development/AutomorphicArithmetic/AutomorphicVarieties/ModularityArgument/ThresholdArguments/DenseThresholdCases/IsSeed.lean}\allowbreak:25--33). \nolinkurl{UpperTheory.V.Z.Sym.W.adjoint_lift_from_eisenstein_of_propagation} (\nolinkurl{Development/AutomorphicArithmetic/AutomorphicVarieties/ModularityArgument/ThresholdArguments/DenseThresholdCases/AdjointLiftFromEisensteinSeedOfHyps.lean}\allowbreak:184--210).}

The field-pair hypothesis identifies $\ker\eta$ with $G_K$. Since $[K:\mathbf Q]=2$, every square belongs to $G_K$, so $\eta^2=1$. Complex conjugation $c$ cannot fix the imaginary quadratic field, so $\eta(c)=-1$. Both $\eta$ and $\theta$ are trivial on $G_L$ by the field pair. The product $\chi'\eta$ is split as prescribed and nontrivial, since $\eta^2=1$ and $\chi'\ne\eta$. These characters that factor through the action on $N$th roots of unity have finite order, hence norm one. Let $k_0$ be the threshold for $\mathrm{PropEGW}$, take $t=B+k_0+2$ and $m=j+(\ell-1)t$. The field-pair parity is $\chi'(c)=(-1)^j$. Since $\eta(c)=-1$ and $\ell-1$ is even, $(\chi'\eta)(c)=(-1)^{j+1}=(-1)^{m+1}$. Choose the rank-two seed $\rho$ with its prescribed PR and unlinked data. Put $\Phi=\epsilon^m\chi'\eta$. On $G_L$, $\Phi\equiv\Psi$ modulo the maximal ideal because $\epsilon^{\ell-1}\equiv1$ on $\ell$-adic units. The field-pair square witness where $\Psi^2\not\equiv1$ therefore also has $\Phi^2\not\equiv1$.

{\footnotesize \nolinkurl{UpperTheory.V.Sd.cong_cyc_pow}, \nolinkurl{finOrder_of_kills}, \nolinkurl{UpperTheory.V.Z.Sym.W.adjoint_lift_from_eisenstein_of_propagation} (\nolinkurl{Development/AutomorphicArithmetic/AutomorphicVarieties/ModularityArgument/ThresholdArguments/DenseThresholdCases/AdjointLiftFromEisensteinSeedOfHyps.lean}\allowbreak:58--84,88--126,184--238).}

\paragraph{Irreducibility survives both operations.}
To see why this one square excludes reducibility on an open subgroup, pass to its open normal core $H$. Ordinary inertia above $\ell$ has an element with infinite-order cyclotomic value. Indeed each finite cyclotomic level is totally ramified, every action on its primitive root lifts to absolute inertia, and the closed sets of lifts prescribing the residues of one $\ell$-adic unit are nested in compact inertia. Use an infinite-order unit and raise its lift to a positive power in $H$. Its ordinary diagonal eigenvalues $a,1$ are distinct since $k-1=m>0$. If $H$ fixes a line, normality implies that all its translates are eigenlines of this element. Global irreducibility supplies a second line. Elements interchanging the two have trace zero, hence $1+\Phi\equiv0$ there by the residual polynomial. For an element preserving the line multiply it by an interchanger and divide the two resulting congruences to get $\Phi\equiv1$. Consequently $\Phi^2\equiv1$ everywhere, contrary to the witness. This proves open irreducibility for $\rho$.

{\footnotesize \nolinkurl{GlobalTheory.CycInertia.exists_mem_Ine_toZModPow}, \nolinkurl{cycZ_Ine_surjective} (\nolinkurl{Development/AbelianVarietyTheory/Proofs/InertiaEigenvalues/CyclotomicCharInertiaSurjective.lean}\allowbreak:37--61,63--96). \nolinkurl{UpperTheory.V.Sd.eisReg} (\nolinkurl{Development/AutomorphicArithmetic/AutomorphicVarieties/IdeleGenerators/HolomorphicIrreducibility/EisReg.lean}\allowbreak:30--56). \nolinkurl{core}, \nolinkurl{UpperTheory.V.Sd.sIabs} (\nolinkurl{Development/AutomorphicArithmetic/AutomorphicVarieties/IdeleGenerators/HolomorphicIrreducibility/LineOfNotIrr.lean}\allowbreak:42--79,81--116). \nolinkurl{UpperTheory.V.Sd.sI} (\nolinkurl{Development/AutomorphicArithmetic/AutomorphicVarieties/IdeleGenerators/HolomorphicIrreducibility/NormOneOfFin.lean}\allowbreak:23--56).}
{\footnotesize The valuation above $\ell$ and the unit of infinite order are supplied by \nolinkurl{eisReg_place}, \path{Development/AutomorphicArithmetic/AutomorphicVarieties/IdeleGenerators/HolomorphicIrreducibility/EisReg.lean:7--21} and \nolinkurl{exists_unit_infinite_order}, \path{Development/AbelianVarietyTheory/Proofs/FoundationProofs/ExistsDecompositionGroupCyclotomicCharacterZpowN.lean:9--31}.}

Set $R=\operatorname{Sym}^2\rho\otimes\det(\rho)^{-1}$. If $R$ on an open subgroup preserves a plane, its linear equation restricts along $(x^2,2xy,y^2)$ to a nonzero quadratic. Its roots in $\mathbb P^1$ form a set of one or two lines permuted by the rank-two action, since Veronese intertwines the actions up to a nonzero determinant factor. The stabilizer of a root in the open subgroup is open: with two distinct roots it is the complement of the inverse image of the other root under the continuous orbit map, and with one root it is the whole subgroup. This would give an invariant line for $\rho$ on an open subgroup. If $R$ instead preserves a line, the nondegenerate symmetric binary-quadratic pairing makes its orthogonal plane invariant and gives the same contradiction. Thus $R$ is openly irreducible. The image of $G_L$ in $G_{\mathbb Q}$ is open and the restriction is a homeomorphism onto that image, so the assertion also holds for $r=R|_{G_L}$.

{\footnotesize \nolinkurl{ad0Irr_factor}, \nolinkurl{ad0Irr_core}, \nolinkurl{UpperTheory.V.Sd.ad0Irr} (\nolinkurl{Development/AutomorphicArithmetic/AutomorphicVarieties/IdeleGenerators/HolomorphicIrreducibility/Ad0IrrMulVec.lean}\allowbreak:7--28,84--171,186--207,213--302). \nolinkurl{UpperTheory.V.Sd.resFacts} (\nolinkurl{Development/AutomorphicArithmetic/AutomorphicVarieties/IdeleGenerators/HolomorphicIrreducibility/ResFacts.lean}\allowbreak:15--47).}

\paragraph{The remaining seed conditions.}
The entries of $R$ are quadratic polynomials in the entries of $\rho$ divided by its invertible determinant, giving continuity. If $P$ is the invariant symmetric form for $R$, then $P R(c)$ is symmetric because $R(c)^2=1$ and $R(c)^t P R(c)=P$. Multiplication shows $R(s)^tP R(c)R(csc^{-1})=P R(c)$, furnishing conjugate polarization. The rank-two trace and unit determinant congruences pass to $\operatorname{tr}(\rho)^2/\det(\rho)-1$, the adjoint trace. Its characteristic polynomial is the one for ratios of the two eigenvalues and $1$; with rank-two residual entries $1,\Phi$ its residual multiset is $\{\Phi,1,\Phi^{-1}\}$. Replace $\Phi$ by the congruent unit $\Psi$ on $G_L$. For an upper-triangular local flag with characters $\kappa_0,\kappa_1$, the induced adjoint flag has diagonal $\kappa_0/\kappa_1,1,\kappa_1/\kappa_0$, hence the three stated inertial weights. On the whole decomposition group each diagonal character is also congruent to the corresponding entry of $(\Psi,1,\Psi^{-1})$: the analogous congruence holds for $\Phi$, and $\Phi\equiv\Psi$ there. This gives the third clause of \nolinkurl{OrdAt} after restriction, by \nolinkurl{ordAt_change} in the cited proof (lines 270--272). The same adjoint characteristic-polynomial calculation for the away type leaves $\chi',1,(\chi')^{-1}$ there since $\eta$ and $\theta$ restrict trivially to $G_L$.

{\footnotesize \nolinkurl{UpperTheory.V.Sd.ad0Facts} (\nolinkurl{Development/AutomorphicArithmetic/AutomorphicVarieties/IdeleGenerators/HolomorphicIrreducibility/QuadraticCharpolyCoeffDifference.lean}\allowbreak:145--206). \nolinkurl{UpperTheory.V.Sd.order_at_adjoint_zero}, \nolinkurl{exists_symmetric_pairing_adjoint_zero}, \nolinkurl{UpperTheory.V.Z.Sym.W.adjoint_lift_from_eisenstein_of_propagation} (\nolinkurl{Development/AutomorphicArithmetic/AutomorphicVarieties/ModularityArgument/ThresholdArguments/DenseThresholdCases/AdjointLiftFromEisensteinSeedOfHyps.lean}\allowbreak:132--167,242--272). \nolinkurl{UpperTheory.V.Sd.resFacts} (\nolinkurl{Development/AutomorphicArithmetic/AutomorphicVarieties/IdeleGenerators/HolomorphicIrreducibility/ResFacts.lean}\allowbreak:26--47).}

The window hypothesis $\mathrm{SeedInW}$ places a scalar twist of $\rho|_{G_L}$ in the domain of $\mathrm{AdAutW}$. The trace-free adjoint of a scalar twist equals the original adjoint since numerator and denominator both acquire the square of the scalar. The automorphy and local witness obtained in that window consequently refer to $r\otimes\epsilon^{-1}$. For the potential-unlinked requirement of the witness, no two of $m,0,-m$ differ by one because $m\ge2$. Equality of a positive power of one diagonal character with the same power of another times cyclotomic would force a nonzero power of cyclotomic to equal $1$ on inertia. Take a rational inertia element with cyclotomic value $1+\ell$ and raise it to a positive power fixing $L$, possible since there are finitely many embeddings of $L$. Its value still has infinite order. This excludes the equality. A common twist multiplies all diagonals by the same character, whose power cancels. This proves the local requirement after the inverse cyclotomic twist and completes the automorphy witness on the same $r$.

{\footnotesize \nolinkurl{UpperTheory.V.Sd.adjoint_zero_twist}, \nolinkurl{UpperTheory.V.Z.Sym.W.adjoint_lift_from_eisenstein_of_propagation} (\nolinkurl{Development/AutomorphicArithmetic/AutomorphicVarieties/ModularityArgument/ThresholdArguments/DenseThresholdCases/AdjointLiftFromEisensteinSeedOfHyps.lean}\allowbreak:168--180,273--281). \nolinkurl{flag_exp}, \nolinkurl{cyc_inf_over}, \nolinkurl{UpperTheory.V.Z.ImaginaryQuadraticEtaCharacter.potOfFlag} (\nolinkurl{Development/AutomorphicArithmetic/AutomorphicVarieties/ModularityArgument/LeviSmoothLocalization/PrincipalSeriesMatching/PotOfFlag.lean}\allowbreak:7--25,27--99). \nolinkurl{pot_tw}, \nolinkurl{UpperTheory.V.Z.Win.flagLaws} (\nolinkurl{Development/AutomorphicArithmetic/AutomorphicVarieties/ModularityArgument/AutomorphyLifting/PolarizedLinkage/FlagLaws.lean}\allowbreak:279--287,398--405). \nolinkurl{UpperTheory.V.Z.Win.Birth.potOrd} (\nolinkurl{Development/AutomorphicArithmetic/AutomorphicVarieties/ModularityArgument/AutomorphyLifting/PrimeIndexConstruction/PotOrd.lean}\allowbreak:7--17). \nolinkurl{UpperTheory.V.Z.Win.Birth.mkCsdW'} (\nolinkurl{Development/AutomorphicArithmetic/AutomorphicVarieties/ModularityArgument/AutomorphyLifting/PrimeIndexConstruction/RankFourAutomorphyOfTensor.lean}\allowbreak:16--17).}

\emph{Complete for the propagation implication.} [Hypothesis, discharged at the application:] its three inputs are exactly the three preceding expert units Z1.1.01--03. The caller supplies them from its published hypotheses and its row assumption. Their independent proofs remain separate review assignments.

{\footnotesize \nolinkurl{UpperTheory.V.Z.Sym.W.unitary_three_automorphy_dense_of_prints_and_row} (\nolinkurl{Development/AutomorphicArithmetic/AutomorphicVarieties/ModularityArgument/ThresholdArguments/DenseThresholdCases/UnitaryThreeInputsOfLiterature.lean}\allowbreak:25--29).}

% blueprint/ch_P.tex lines 12654 through 12655; r6:review:Z1.1.04
\paragraph{Z1.1.04 and the two closing Kato implications.}\label{r6:review:Z1.1.04} \emph{Original argument on p.~\pageref{r5:Z1.1.04}.} 
Z1.1.04 retains complete status for the implication from the three inputs on a common seed; the character kernel, parity and ordinary diagonal congruences are explicit in the revised argument (\nolinkurl{adjoint_lift_from_eisenstein_of_propagation}, \path{Development/AutomorphicArithmetic/AutomorphicVarieties/ModularityArgument/ThresholdArguments/DenseThresholdCases/AdjointLiftFromEisensteinSeedOfHyps.lean:184--281}).

